% #############################################################################
% RESUMO em Português
% !TEX root = ../main.tex
% #############################################################################
% use \noindent in first paragraph


\noindent Veículos alimentados a energia solar, em particular embarcações, necessitam de uma conversão de energia eficaz em condições climáticas de rápida variação. O carregador \ac{mppt} desempenha um papel importante nesta conversão, garantindo que os painéis operem no seu \ac{mpp}. Este trabalho centra-se no projeto, simulação e ensaio de um carregador \ac{mppt} de resposta rápida para o \ac{tsb}.

O protótipo foi projetado para carregar uma bateria de iões de lítio com tensão nominal de 43,2~V. O carregador proposto baseia-se num conversor boost síncrono que opera a 100~kHz, controlado por um \ac{mcu}.

Foram analisadas diferentes técnicas de \ac{mppt} de modo a identificar uma solução adequada às condições de operação dinâmicas da embarcação. O compromisso entre a velocidade de seguimento, precisão e complexidade levou à implementação de um algoritmo \ac{peo}. O conversor e o modelo fotovoltaico foram simulados no LTspice, e o \ac{mppt} algoritmo no MATLAB/Simulink. Foi projetada, fabricada e soldada manualmente uma \ac{pcb} dedicada. O firmware no \ac{mcu} implementa \ac{peo} com carregamento \acs{cc}-\acs{cv}, comunicação \ac{can} e uma aplicação para configuração.

Os ensaios de laboratório deram uma eficiência máxima de 96,84\%, com um tempo de seguimento tipicamente inferior a 100~ms e de 200~ms no caso mais exigente, abaixo dos alvos de 0,5~s e 1~s. O ruído conduzido e um ensaio térmico de 3~h permaneceram dentro de limites seguros. O carregador seguiu o perfil \acs{cc}-\acs{cv} em laboratório e num ensaio de campo em terra. Com um custo de cerca de 125~€, a placa constitui um carregador \ac{mppt} dedicado, configurável e de baixo custo para o \ac{tsb}.

\vspace{0.5cm}
